%% Paper 2 SKELETON — ICRA/IROS (IEEEtran conference).
%% This is a structured skeleton: abstract is a working draft; sections are stubs with
%% the arguments to make. Fill experiments once the policy + baseline are trained.
\documentclass[conference]{IEEEtran}
\IEEEoverridecommandlockouts
\usepackage{cite}
\usepackage{amsmath,amssymb,amsfonts}
\usepackage{algorithmic}
\usepackage{graphicx}
\usepackage{textcomp}
\usepackage{xcolor}
\usepackage{url}
\newcommand{\todo}[1]{\textcolor{red}{[TODO: #1]}}

\begin{document}

\title{Robust Detumbling of Tumbling Debris with Unknown Inertia using a
Free-Floating Space Manipulator and Domain-Randomized Reinforcement Learning}

\author{\IEEEauthorblockN{Karam Khasawneh \textit{et al.}}
\IEEEauthorblockA{\textit{Dept.\ of Robotics and Artificial Intelligence}\\
\textit{Jadara University}, Irbid, Jordan\\
KaramQ5@ieee.org}}

\maketitle

\begin{abstract}
Capturing and detumbling non-cooperative space debris is a defining challenge for on-orbit
servicing: the target is tumbling, its mass and inertia are unknown, and the chaser is a
\emph{free-floating} platform whose attitude is disturbed by every arm motion. Prior
reinforcement-learning (RL) work for space manipulators concentrates on the pre-contact phase
(reaching and tracking a grasp point), while post-grasp detumbling is handled by model-based
controllers that assume known target dynamics. We address the gap with an RL controller for the
\emph{post-grasp detumbling} phase that is explicitly trained to be \emph{robust to unknown target
inertia} via domain randomization over the target mass, inertia tensor, and initial angular
velocity. We evaluate in a MuJoCo free-floating-base simulation with full contact dynamics against
an impedance / active-resistance baseline, reporting detumbling time, contact force, and induced
base-attitude disturbance. \todo{headline result: learned policy holds success across $0.5\times$--$2\times$
nominal inertia where the model-based baseline degrades by $Y\%$.}
\end{abstract}

\begin{IEEEkeywords}
Space robotics, on-orbit servicing, active debris removal, reinforcement learning,
free-floating manipulator, contact control, domain randomization.
\end{IEEEkeywords}

\section{Introduction}
\todo{(1) debris + non-cooperative capture motivation; (2) the three sub-problems: reach, grasp,
detumble; (3) the gap: RL stops at grasp, model-based detumbling assumes known inertia;
(4) our contribution bullets; (5) contributions list.}

Contributions:
\begin{itemize}
\item An RL formulation of \emph{post-grasp detumbling} on a free-floating base with contact dynamics.
\item A domain-randomization scheme over target inertia that yields robustness to \emph{unknown} targets.
\item A benchmark against a model-based detumbling baseline on detumble time, contact force, and
      base-attitude disturbance, with an ablation isolating the randomization.
\end{itemize}

\section{Related Work}
\todo{Three clusters (see literature\_review.md). A: high-level ADR sequencing (not us).
B: pre-contact RL guidance/tracking on free-floating bases \cite{guidance2510,wei2025robotica}
-- stops at grasp. C: contact/detumbling, mostly model-based \cite{activeresist2025}; RL soft-capture
with tactile \cite{softcapture2409} lacks unknown-inertia robustness and model-based baselines.
State precisely what is new here.}

\section{Problem Formulation}
\subsection{Free-floating dynamics}
\todo{Equations of motion for base+arm with momentum conservation (zero external wrench);
the coupling that makes arm motion disturb base attitude. Contact model (soft, friction).}

\subsection{MDP}
State $s\in\mathbb{R}^{32}$: base pose+twist, joint positions/velocities, target relative
pose and angular velocity, end-effector-to-grasp vector. Action $a\in[-1,1]^3$: normalized joint
torques. Reward penalizes target spin, end-effector distance, induced base angular velocity, and
control effort, with a terminal bonus when the target angular velocity falls below a threshold:
\begin{equation}
r = -w_\omega\lVert\boldsymbol{\omega}_t\rVert - w_d d_{\text{eg}}
    - w_b\lVert\boldsymbol{\omega}_b\rVert - w_u\lVert a\rVert^2 + r_{\text{term}}.
\end{equation}
\todo{This matches \texttt{robotic\_capture/sim/free\_flyer\_env.py}; lift weights once tuned.}

\section{Method}
\todo{Policy/algorithm (e.g.\ SAC/PPO from SB3); domain randomization distribution over
target mass and inertia ($0.5\times$--$2\times$) and initial tumble; training curriculum; network.}

\section{Baselines}
\todo{Re-implement an impedance / active-resistance detumbling controller \cite{activeresist2025}
tuned to nominal inertia. This is the comparison that makes the robustness claim meaningful.}

\section{Experiments}

We evaluate the learned policy (SAC trained on 120k steps with domain randomization over
target mass and inertia; 3 seeds for statistical rigor) and the baseline (Active Resistance
Control, ``ARC''; force-feedback inertia-agnostic) on a paired sweep grid: inertia scale
$\{0.5, 1.0, 1.5, 2.0\} \times$ spin magnitude $\{0.2, 0.3, 0.4\} \text{ rad/s} \times$ seed
$\{0, 1, 2, 3, 4\}$. Each condition is run to task termination or max 500 steps. We report
(1) success rate (spin below 0.05 rad/s threshold), (2) base-attitude disturbance (induced
base angular velocity; proxy for chaser ACS fuel), (3) spin reduction, and (4) detumble time.
Bootstrap 95\% confidence intervals (10k resamples) bracket the means.

\subsection{Robustness across inertia: base-attitude disturbance}

The headline result (Fig.~\ref{fig:base_disturbance}) is \emph{induced base-attitude
disturbance} as a function of target inertia. The RL policy trained with domain randomization
(RL-DR) keeps base disturbance flat and low (\textbf{0.016--0.026 rad/s}) across the full
inertia sweep, whereas the model-based baseline (ARC) rises steeply (\textbf{0.034 → 0.285 rad/s},
an 8--19\(\times\) increase). This disturbance quantifies the momentum transfer back to the
chaser's attitude controller: ARC's force-tracking ignores this cost and accumulates higher
base spin at targets with larger inertia. The RL policy explicitly minimizes base angular
velocity in the reward, hence the robustness to unknown target mass.

\begin{figure}[t]
\centering
\includegraphics[width=0.9\linewidth]{../../results/fig_base_disturbance.pdf}
\caption{Induced base-attitude disturbance vs.\ target inertia scale. Three lines: RL-DR
(domain-randomized policy, pooled over 3 training seeds), ARC (model-based baseline),
and zero control. RL-DR is flat and robust; ARC rises steeply with inertia.}
\label{fig:base_disturbance}
\end{figure}

\begin{table}[t]
\centering
\caption{Induced base-attitude disturbance (rad/s), mean [95\% bootstrap CI]. RL-DR pooled
over 3 seeds; ARC nominal tuning; Zero = no control.}
\label{tab:base_disturbance}
\begin{tabular}{c|ccc}
\hline
Inertia & RL-DR & ARC & Zero \\
\hline
$0.5\times$ & \textbf{0.026} [0.022, 0.030] & 0.034 [0.019, 0.049] & 0.000 \\
$1.0\times$ & \textbf{0.016} [0.012, 0.020] & 0.121 [0.091, 0.153] & 0.000 \\
$1.5\times$ & \textbf{0.015} [0.013, 0.017] & 0.226 [0.171, 0.284] & 0.000 \\
$2.0\times$ & \textbf{0.020} [0.017, 0.023] & 0.285 [0.228, 0.338] & 0.000 \\
\hline
\end{tabular}
\end{table}

Success rate is 100\% for RL-DR across all scales; ARC degrades to 67\% at $1.5\times$ and
$2.0\times$ (final spin $\sim 0.07$--$0.13$ rad/s, above the 0.05 rad/s detumbling threshold).
Spin reduction (percentage target spin decrease) is comparable between RL-DR and ARC
(82--96\%), consistent with both methods achieving the primary task; the RL advantage manifests
in the \emph{secondary} metric: base disturbance (ACS fuel cost).

\subsection{Ablation: domain randomization}

To isolate the robustness contribution of domain randomization, we trained the same SAC policy
without randomization (fixed nominal inertia, scale $1.0\times$) for 120k steps. Tested on the
same sweep, the no-DR policy still generalizes reasonably to heavy targets (100\% success at
$1.5\times$--$2.0\times$), but loses robustness at the light extreme: success drops to 67\% at
$0.5\times$ (RL-DR: 100\%). This trade-off is expected: no-DR specializes to nominal inertia
(lower base disturbance at $1.0\times$: 0.011 vs.\ 0.016 rad/s) but pays a robustness penalty
at the distribution tails. The result validates domain randomization as the mechanism enabling
inertia robustness.

\subsection{Sensing and control effort}

RL-DR uses kinematics only (32-D observation: base pose/twist, joint positions/velocities, target
relative motion) and achieves lower base disturbance. ARC requires a contact force/torque sensor
(force feedback in Eq.~5 of the baseline paper) and still incurs higher base disturbance. This
represents a \emph{sensing-reduction advantage}: the learned policy operates blind to contact
force and target inertia yet outperforms the model-based controller on the chaser's attitude
cost. Control effort (integral of $\lVert a \rVert^2$) is comparable for RL-DR and ARC;
RL-DR does not sacrifice actuator usage to achieve lower base disturbance.

\section{Limitations and Sim-to-Real}
\todo{Contact-model fidelity; planar 3-DOF scaffold vs.\ 6-DOF arm; no hardware (state plainly);
products of inertia; sensor noise. Name future high-fidelity / hardware validation.}

\section{Conclusion}
\todo{Restate the robustness-to-unknown-inertia result and its OOS relevance.}

\begin{thebibliography}{9}
\bibitem{guidance2510} ``Safe Obstacle-Free Guidance of Space Manipulators in Debris Removal
Missions via Deep Reinforcement Learning,'' arXiv:2510.06566. \todo{full authors/venue}
\bibitem{wei2025robotica} Y.\ Wei, X.\ Bai, H.\ Lu, ``Trajectory planning of free-floating space
robot for non-cooperative tumbling target capture based on deep reinforcement learning,''
\textit{Robotica}, 2025. \todo{vol/pp/doi}
\bibitem{activeresist2025} ``Active resistance control: a contact control method for detumbling
non-cooperative spacecraft by a robotic arm,'' \textit{Nonlinear Dynamics}, 2025,
doi:10.1007/s11071-025-11381-z.
\bibitem{softcapture2409} ``Improving Soft-Capture Phase Success in Space Debris Removal Missions:
Leveraging Deep Reinforcement Learning and Tactile Feedback,'' arXiv:2409.12273.
\bibitem{review2024} ``Review of Autonomous Space Robotic Manipulators for On-Orbit Servicing and
Active Debris Removal,'' \textit{Space: Science \& Technology}, 2024, doi:10.34133/space.0291.
\bibitem{domrand} X.\ B.\ Peng \textit{et al.}, ``Sim-to-Real Transfer of Robotic Control with
Dynamics Randomization,'' arXiv:1710.06537.
\end{thebibliography}
\end{document}
